\documentclass[a4paper,11pt]{ltjsarticle}
\usepackage[margin=25mm]{geometry}
\usepackage{amsmath,amssymb,amsthm,mathtools}
\usepackage[unicode,hidelinks]{hyperref}
\newtheorem{theorem}{定理}[section]
\newtheorem{lemma}[theorem]{補題}
\newtheorem{proposition}[theorem]{命題}
\newtheorem{corollary}[theorem]{系}
\theoremstyle{remark}
\newtheorem{remark}[theorem]{注意}
\renewcommand{\proofname}{証明}
\newcommand{\Perf}{\operatorname{Perf}}
\newcommand{\RHom}{\operatorname{RHom}}
\newcommand{\Cone}{\operatorname{Cone}}
\newcommand{\thick}{\operatorname{thick}}
\newcommand{\otL}{\otimes^{\mathbf L}}
\newcommand{\Spec}{\operatorname{Spec}}
\newcommand{\HH}{\operatorname{HH}}
\title{相対残余dg圏の有限対角判定とgeneric spreading\\\large 可換復元の相対smoothnessへ向けた研究ノート}
\author{nc-cubic：主研究ノート}
\date{2026年11月10日・v1}
\begin{document}
\maketitle
\begin{abstract}
有限型基底上のproperなdg代数について、対角のperfect性を有限bar段階のretractとして表す判定、generic fiberでのperfect性を有限局所化へ降下する補題、Artin底でのperfectness検出の適用限界を整理する。
原資料は\texttt{notes/pdf-only/finite\_smoothness\_notes.pdf}の第15--22節、とくに第18--21節（pp.~23--25）である。既存の第23節の反例は\texttt{notes/source\_recovery/dg\_perfectness\_counterexamples.tex}にある。本稿はそれらの現在有効な証明を可換復元の残余対角へ適用しやすい形で整理したもので、原PDF全体のTeX置換ではない。相対smoothnessの新しい無条件定理を主張しない。
\end{abstract}

\section{設定と区別すべき三つのperfect性}
$R$を可換環、$\Gamma$を単位的な$R$-dg代数とし、
\[
 \Gamma^e=\Gamma^{\mathrm{op}}\otL_R\Gamma,\qquad
 \Delta_\Gamma=\Gamma\in D(\Gamma^e)
\]
とおく。$\Gamma\in\Perf(R)$をproper性の入力とする。
ここで\emph{相対smoothness}とは$\Delta_\Gamma\in\Perf(\Gamma^e)$をいう。
$\Gamma$のunderlying $R$-複体のperfect性、$\Gamma\otL_Rk$の対角perfect性、$\Delta_\Gamma$の$\Gamma^e$-perfect性は、論理的に異なる条件である。

非楕円型の節点・尖点量子化では$R=\mathbb C[r,s]_{(r,s)}$、$\Gamma_R=\operatorname{REnd}(G)$が実在するproper残余dg模型である。特殊ファイバー$\Gamma_R\otL_Rk$は滑らかなcubic fivefoldのKuznetsov成分を表しsmoothであるが、この事実だけで$\Gamma_R$の相対smoothnessを結論したことにはならない。

\section{有限bar段階による対角判定}
$\operatorname{Bar}_\bullet(\Gamma/R)$を、導来テンソルで形成する相対二側bar解消とする。必要ならcofibrantな模型を取る。$P_d$をその$d$-skeletonの有限部分実現とし、増大写像を$\epsilon_d:P_d\to\Delta_\Gamma$と書く。

\begin{proposition}[有限段階のretract判定]\label{prop:bar}
$\Gamma\in\Perf(R)$と仮定する。このとき$P_d\in\Perf(\Gamma^e)$で、
\[
 \operatorname{hocolim}_{d}P_d\simeq\Delta_\Gamma.
\]
従って次の条件は同値である。
\begin{enumerate}
\item $\Gamma$が$R$上homologically smoothである。
\item ある$d\ge0$に対し$D(\Gamma^e)$で$\epsilon_d:P_d\to\Delta_\Gamma$が右逆
$s_d:\Delta_\Gamma\to P_d$を持つ。
\end{enumerate}
\end{proposition}
\begin{proof}
barの各simplicial項は、導来双加群として
\[
 \Gamma\otL_R(\Gamma)^{\otL_R d}\otL_R\Gamma
 \simeq \Gamma^e\otL_R (\Gamma)^{\otL_R d}
\]
という形である。$\Gamma\in\Perf(R)$なので有限テンソル冪は$R$-perfectであり、その$\Gamma^e$への拡大は$\Gamma^e$-perfectである。有限skeletonはこれらの有限個のshift、cone、retractで作られるから$P_d$はperfectである。二側barのaugmentationは実現後に擬同型となる。

右逆が存在すれば$\Delta_\Gamma$は$P_d$のretractなのでperfect。逆に$\Delta_\Gamma$がperfectなら$D(\Gamma^e)$でcompactである。したがって
\[
 \operatorname{Hom}_{D(\Gamma^e)}
 (\Delta_\Gamma,\operatorname{hocolim}_d P_d)
 \cong \varinjlim_d\operatorname{Hom}_{D(\Gamma^e)}(\Delta_\Gamma,P_d)
\]
を用いると、恒等射はある有限段階を経由する。この因子化は$\epsilon_d$の右逆を与える。
\end{proof}

\begin{remark}[有限段階の選択は解決ではない]\label{rem:lifting}
特殊ファイバーの対角がsmoothなら、ある$d$で$\epsilon_{d,k}$に右逆がある。しかし、その右逆やそのhomotopyが$R$上に\emph{実際に}持ち上がることは別の課題である。必要なのは$\Gamma_R^e$上の有限データ
\[
 s_d:\Gamma_R\longrightarrow P_{d,R},\quad
 H:\epsilon_{d,R}s_d\simeq\mathrm{id}_{\Gamma_R}
\]
（またはこれと同値な有限perfect複体と増大射）である。各Artin切断で異なる有限段数のretractが存在することは、一様な$d$や代数的な右逆を与えない。
\end{remark}

\section{局所Artin底での検出と完備底の障害}
\begin{lemma}[Artin底でのperfectness検出]\label{lem:artin}
$R$が局所Artin環で剰余体$k$を持つとする。任意の$R$-dg代数$\Lambda$と$M\in D(\Lambda)$について、
\[
 M\in\Perf(\Lambda)
 \quad\Longleftrightarrow\quad
 M\otL_Rk\in\Perf(\Lambda\otL_Rk).
\]
特に$\Lambda\otL_Rk$がsmoothなら$\Lambda$は相対smoothである。
\end{lemma}
\begin{proof}
逆向きを述べればよい。$R$は$k$を組成因子とする有限filtrationを持つため$R\in\thick_{D(R)}(k)$である。$M\otL_Rk$のcompactnessにより、$X\mapsto\RHom_\Lambda(M,-\otL_RX)$が任意の直和を保つという条件を満たす$X$は$k$を含むthick部分圏をなす。ゆえに$R$も含み、$M$は$\Lambda$-compact。smoothnessの主張は$\Lambda^e$上の対角に適用する。
\end{proof}

\begin{remark}[完備底への誤った外挿]
$R=k[[h]]$では一般に$R\notin\thick_{D(R)}(k)$であり、上の議論は使えない。全ての$R/(h^n)$でperfectであっても、$R$上でperfectとは限らない。\texttt{notes/source\_recovery/dg\_perfectness\_counterexamples.tex}の明示例はその失敗を示すが、特殊ファイバーの代数自身はsmoothでないため、滑らかな特殊Kuznetsov成分からの相対対角持上げの直接反例ではない。
\end{remark}

\section{generic fiberのperfect性のspread}
\begin{proposition}[genericなperfectnessの局所化]\label{prop:spread}
$R$をnoetherian整域、$K=\operatorname{Frac}(R)$、$\Lambda$をunderlying $R$-perfectなdg代数、$M\in D(\Lambda)$をunderlying $R$-perfectなdg加群とする。$M_K\in\Perf(\Lambda_K)$なら、ある$0\neq s\in R$が存在して
\[
 M\otL_RR[s^{-1}]\in\Perf(\Lambda\otL_RR[s^{-1}]).
\]
\end{proposition}
\begin{proof}
$K=\varinjlim_{s\neq0}R[s^{-1}]$である。perfect dg加群のfiltered-colimit descent（有限cellとretractの表示、またはToën--VaquiéのLemma~2.10に基づく議論）により、ある局所化$R[s^{-1}]$上のperfect $\Lambda_s$-加群$Q_s$で$Q_s\otL_{R_s}K\simeq M_K$となるものを得る。

$Q_s$のcompactnessにより、generic上の同型は必要に応じてさらに局所化した後、射$Q_s\to M_s$で表せる。$Q_s$は$\Lambda_s$-perfect、$\Lambda_s$は底上perfectなので$Q_s$は$R_s$-perfectである。従って$C=\Cone(Q_s\to M_s)$は$R_s$-perfectであり、genericで消滅する。その有限個の有限生成cohomologyを殺す非零元を反転すると$C=0$となる。よって局所化後$M_s\simeq Q_s$となる。
\end{proof}

\begin{corollary}[相対smoothnessのgeneric spreading]\label{cor:spread}
$R$をnoetherian整域とし、$\Gamma\in\Perf(R)$をdg代数とする。$\Gamma_K$が$K$上smoothなら、ある$0\neq s\in R$で$\Gamma_s$は$R_s$上smooth properである。
\end{corollary}
\begin{proof}
命題\ref{prop:spread}を$\Lambda=\Gamma^e$、$M=\Delta_\Gamma$に適用する。どちらもunderlying $R$-perfectである。
\end{proof}

\begin{remark}[何を証明していないか]\label{rem:not-open}
系\ref{cor:spread}はgeneric fiberがすでにsmoothと分かっているときの結果である。特殊ファイバーがsmoothという情報だけではgeneric smoothnessは得られない。よってこの系を$\Gamma_R$の相対smoothness (SM) の証明にそのまま用いることはできない。
\end{remark}

\begin{corollary}[一変数での条件付きfiberwise結論]\label{cor:onevariable}
$R=k[h,\Delta(h)^{-1}]$、$\Delta(0)\neq0$とする。$\Gamma\in\Perf(R)$で必要なdg底変更が成立すると仮定する。$\Gamma_0$と$\Gamma\otL_Rk((h))$がそれぞれの体上smoothなら、$q(0)\neq0$を満たす多項式$q(h)$が存在し、$\Delta(h)q(h)\neq0$である全体値ファイバー$\Gamma_\lambda$はsmooth properとなる。
\end{corollary}
\begin{proof}
$k(h)\subset k((h))$の体拡大に対するfaithfully flatな下降により、generic $k(h)$上で対角がperfectとなる。系\ref{cor:spread}によってある$D(s)$上で相対smoothとなる。$s=h^a q(h)$（$q(0)\neq0$）と書けば、$h\neq0$では$D(q)$の全点が$D(s)$に含まれ、$h=0$は別途仮定でsmooth。これは\emph{全体のファイバーの}結論であって、$D(q)$全体で相対対角がperfectという主張ではない。
\end{proof}

\section{非楕円型cubic fivefoldに残る有限作業}
節点型・尖点型の二パラメータ族のための実在するproper残余dg代数を$\Gamma_R$とする。既存原稿\texttt{papers/sections/local\_projective\_model.tex}はその構成と線束Homの底変更を扱い、\texttt{papers/sections/nonelliptic\_quantization.tex}は量子化と中心三次商の具体式を扱う。いずれも今回のノートの相対対角右逆を作ったものではない。

優先的に構成すべきデータは、$\Gamma_R$の有限のdg（または$A_\infty$）表示を選び、ある$d$に対しbarの有限段階$P_{d,R}$、右逆$s_d$、増大後の恒等homotopyを\emph{有限個の整合式として}提示することである。$\Gamma_R$のunderlying cohomology階数（fivefoldの場合、線束生成dg代数の0次と5次にそれぞれ2838、45）や\(\HH^3(\operatorname{Ku}(X))=0\)だけでは、このデータは得られない。

\paragraph{出典と検証範囲}
本稿は\texttt{notes/pdf-only/finite\_smoothness\_notes.pdf}（2026年11月1日版）第18--21節の補題・命題・証明を現在の問題へ組み直したもの。原本の他の節は残るためPDFのGit追跡解除を認めない。Toën--Vaquié, \emph{Moduli of objects in dg-categories}, Ann. Sci. ENS \textbf{40} (2007), 387--444, Lemmas 2.8 and 2.10, \url{https://arxiv.org/abs/math/0503269} を参照。参考文献の適用と導来bar模型の細部は最終公表前に改めて監査する。
\end{document}
